\chapter{Formula Sheets and Background}
\label{ch:appendix}

\section{Vector identities}

Let $f,g$ be scalar fields and $\vec F,\vec G$ vector fields (smooth enough). Proofs are in Chapter~\ref{ch:fields} (Theorem~\ref{thm:identities}).

\begin{center}
\renewcommand{\arraystretch}{1.5}
\begin{tabular}{ll}
\toprule
\textbf{Identity} & \textbf{Name / note} \\
\midrule
$\curl(\nabla f)=\vec0$ & gradients have no curl \\
$\divg(\curl\vec F)=0$ & curls have no divergence \\
$\nabla(fg)=f\nabla g+g\nabla f$ & product rule \\
$\divg(f\vec F)=\nabla f\cdot\vec F+f\divg\vec F$ & \\
$\curl(f\vec F)=\nabla f\times\vec F+f\curl\vec F$ & \\
$\divg(\vec F\times\vec G)=\vec G\cdot\curl\vec F-\vec F\cdot\curl\vec G$ & \\
$\curl\curl\vec F=\nabla(\divg\vec F)-\lap\vec F$ & vector Laplacian \\
$\lap f=\divg\nabla f$ & Laplacian \\
\bottomrule
\end{tabular}
\end{center}

\section{Products of vectors}

\begin{center}
\renewcommand{\arraystretch}{1.5}
\begin{tabular}{ll}
\toprule
$\vec a\cdot\vec b=\norm{\vec a}\norm{\vec b}\cos\theta$ & $\norm{\vec a\times\vec b}=\norm{\vec a}\norm{\vec b}\sin\theta$ \\
$\vec a\times\vec b=-\vec b\times\vec a$ & $\vec a\cdot(\vec b\times\vec c)=\det[\vec a\ \vec b\ \vec c]$ \\
$\vec a\times(\vec b\times\vec c)=\vec b(\vec a\cdot\vec c)-\vec c(\vec a\cdot\vec b)$ & $(\vec a\times\vec b)\cdot(\vec c\times\vec d)=(\vec a\cdot\vec c)(\vec b\cdot\vec d)-(\vec a\cdot\vec d)(\vec b\cdot\vec c)$ \\
$\abs{\vec a\cdot\vec b}\le\norm{\vec a}\norm{\vec b}$ & $\norm{\vec a+\vec b}\le\norm{\vec a}+\norm{\vec b}$ \\
\bottomrule
\end{tabular}
\end{center}

\section{Operators in orthogonal coordinates}

For orthogonal coordinates $(u_1,u_2,u_3)$ with scale factors $h_i$ and unit vectors $\ehat_i$ (Chapter~\ref{ch:curvilinear}):
\begin{align*}
\nabla f&=\sum_i\frac1{h_i}\pd f{u_i}\ehat_i,\\
\divg\vec F&=\frac1{h_1h_2h_3}\Bigl[\pd{}{u_1}(h_2h_3F_1)+\pd{}{u_2}(h_3h_1F_2)+\pd{}{u_3}(h_1h_2F_3)\Bigr],\\
\lap f&=\frac1{h_1h_2h_3}\sum_i\pd{}{u_i}\Bigl(\frac{h_1h_2h_3}{h_i^2}\pd f{u_i}\Bigr).
\end{align*}

\begin{center}
\renewcommand{\arraystretch}{1.5}
\begin{tabular}{llll}
\toprule
 & \textbf{Cartesian} & \textbf{Cylindrical $(r,\theta,z)$} & \textbf{Spherical $(\rho,\varphi,\theta)$} \\
\midrule
$h_i$ & $1,1,1$ & $1,\ r,\ 1$ & $1,\ \rho,\ \rho\sin\varphi$ \\
$\dd V$ & $\dd x\dd y\dd z$ & $r\,\dd r\dd\theta\dd z$ & $\rho^2\sin\varphi\,\dd\rho\dd\varphi\dd\theta$ \\
$\lap f$ & $f_{xx}+f_{yy}+f_{zz}$ & $\frac1r(rf_r)_r+\frac1{r^2}f_{\theta\theta}+f_{zz}$ & $\frac1{\rho^2}(\rho^2f_\rho)_\rho+\frac1{\rho^2\sin\varphi}(\sin\varphi f_\varphi)_\varphi+\frac1{\rho^2\sin^2\varphi}f_{\theta\theta}$ \\
\bottomrule
\end{tabular}
\end{center}

\section{Integral theorems}

\begin{center}
\renewcommand{\arraystretch}{1.6}
\begin{tabular}{ll}
\toprule
\textbf{Theorem} & \textbf{Statement} \\
\midrule
Line integrals & $\displaystyle\int_C\nabla f\cdot\dd\vec r=f(B)-f(A)$ \\
Green & $\displaystyle\oint_{\partial D}P\,\dd x+Q\,\dd y=\iint_D(Q_x-P_y)\,\dd A$ \\
Stokes & $\displaystyle\oint_{\partial S}\vec F\cdot\dd\vec r=\iint_S\curl\vec F\cdot\hat{\vec n}\,\dd S$ \\
Gauss & $\displaystyle\iiint_E\divg\vec F\,\dd V=\oiint_{\partial E}\vec F\cdot\hat{\vec n}\,\dd S$ \\
General & $\displaystyle\int_\Omega\dd\omega=\int_{\partial\Omega}\omega$ \\
Green I & $\displaystyle\iiint_E(u\lap v+\nabla u\cdot\nabla v)\,\dd V=\oiint_{\partial E}u\,\partial_nv\,\dd S$ \\
Green II & $\displaystyle\iiint_E(u\lap v-v\lap u)\,\dd V=\oiint_{\partial E}(u\,\partial_nv-v\,\partial_nu)\,\dd S$ \\
\bottomrule
\end{tabular}
\end{center}

\section{Standard parametrizations and elements}

\begin{center}
\renewcommand{\arraystretch}{1.5}
\begin{tabular}{lll}
\toprule
\textbf{Object} & \textbf{Parametrization} & \textbf{Element} \\
\midrule
Circle, radius $R$ & $(R\cos t,R\sin t)$ & $\dd s=R\,\dd t$ \\
Helix & $(a\cos t,a\sin t,bt)$ & $\dd s=\sqrt{a^2+b^2}\,\dd t$ \\
Graph $z=f(x,y)$ & $(x,y,f)$ & $\dd S=\sqrt{1+f_x^2+f_y^2}\,\dd A$ \\
Sphere, radius $R$ & $R(\sin\varphi\cos\theta,\sin\varphi\sin\theta,\cos\varphi)$ & $\dd S=R^2\sin\varphi\,\dd\varphi\dd\theta$ \\
Cylinder, radius $R$ & $(R\cos\theta,R\sin\theta,z)$ & $\dd S=R\,\dd\theta\dd z$ \\
Cone, half-angle $\alpha$ & $(s\sin\alpha\cos\theta,s\sin\alpha\sin\theta,s\cos\alpha)$ & $\dd S=s\sin\alpha\,\dd s\,\dd\theta$ \\
\bottomrule
\end{tabular}
\end{center}

\section{Background facts used in the book}

We list the results from other courses that we used without proof.

\begin{description}
  \item[Fundamental theorem of calculus.] If $f$ is continuous on $[a,b]$ then $\frac{\dd}{\dd x}\int_a^xf=f(x)$, and $\int_a^bf'=f(b)-f(a)$ for $f\in C^1$.
  \item[Mean value theorem.] If $f$ is continuous on $[a,b]$ and differentiable inside, then $f(b)-f(a)=f'(c)(b-a)$ for some $c\in(a,b)$.
  \item[Taylor's theorem (one variable).] $f(x+h)=\sum_{j<k}\frac{f^{(j)}(x)}{j!}h^j+\frac{f^{(k)}(\xi)}{k!}h^k$ for some $\xi$ between $x$ and $x+h$.
  \item[Compactness.] A continuous function on a closed bounded set of $\R^n$ is bounded, attains its maximum and minimum, and is uniformly continuous.
  \item[Determinants.] $\det(AB)=\det A\det B$; $\det A\neq0$ iff $A$ is invertible; $\abs{\det A}$ is the volume scaling factor; $\det$ is multilinear and alternating in the columns; $\frac{\dd}{\dd t}\det A=\det A\,\tr(A^{-1}\dot A)$.
  \item[Spectral theorem.] A real symmetric matrix has an orthonormal basis of eigenvectors with real eigenvalues. It is positive definite iff all eigenvalues are positive.
  \item[Banach fixed point theorem.] If $T:X\to X$ is a contraction ($d(Tx,Ty)\le q\,d(x,y)$, $q<1$) on a complete metric space, then $T$ has a unique fixed point, found as the limit of $x,Tx,T^2x,\dots$ (used in the inverse function theorem, Chapter~\ref{ch:diff}).
  \item[Differentiation under the integral sign.] If $f$ and $\partial_tf$ are continuous on $[a,b]\times[c,d]$, then $\frac{\dd}{\dd t}\int_a^bf(x,t)\dd x=\int_a^b\partial_tf\,\dd x$ (Theorem~\ref{thm:leibniz}).
  \item[Fubini.] For $f$ continuous on a rectangle, iterated integrals in either order give the double integral (Theorem~\ref{thm:fubini}).
\end{description}

\section{Greek alphabet and symbols}

\begin{center}
\begin{tabular}{llll}
\toprule
$\alpha$ alpha & $\beta$ beta & $\gamma$ gamma & $\delta$ delta \\
$\varepsilon$ epsilon & $\theta$ theta & $\kappa$ kappa & $\lambda$ lambda \\
$\mu$ mu & $\rho$ rho & $\sigma$ sigma & $\phi,\varphi$ phi \\
$\psi$ psi & $\omega$ omega & $\Phi$ Phi & $\Omega$ Omega \\
\bottomrule
\end{tabular}
\end{center}

\section{Where to go next}

\begin{itemize}
  \item \emph{Manifolds and forms:} M.\ Spivak, \emph{Calculus on Manifolds}; V.\ Guillemin and A.\ Pollack, \emph{Differential Topology}; J.\ Munkres, \emph{Analysis on Manifolds}.
  \item \emph{Physics:} D.\ Griffiths, \emph{Introduction to Electrodynamics}; J.\ D.\ Jackson, \emph{Classical Electrodynamics}.
  \item \emph{Rigorous multivariable analysis:} W.\ Rudin, \emph{Principles of Mathematical Analysis} (chapters 9--10).
  \item \emph{Optimization and automatic differentiation:} S.\ Boyd and L.\ Vandenberghe, \emph{Convex Optimization}; A.\ Griewank and A.\ Walther, \emph{Evaluating Derivatives}.
  \item \emph{Partial differential equations:} L.\ C.\ Evans, \emph{Partial Differential Equations}.
\end{itemize}
